% Chapter 6 introduction.

\section{Introduction}
\label{p:sec:introduction}

A fitted release rate is not a property of a pathogen. In the standard
within-host model an infected cell sheds virions at a constant rate $p$ and
is removed at a constant rate $d_{\Icell}$, and both numbers are read off
from data as though they characterised the organism. For a budding virus
that is a defensible caricature: release and death really are separate,
independent events, and two constants do capture the cell. For a
lytic pathogen it is wrong in mechanism. Such a pathogen releases nothing
while it grows inside the cell, and then releases everything at once, in the
event that kills the cell. Release and death are not independent; they are
the same event. \yp{} in the macrophage, \ft{} in the macrophage, the lytic
phage in its bacterium --- all are of that kind, and none is described by a
constant drip. Nor, once the mechanism is put back in, does $p$ survive as a
constant of the organism: it becomes a property of the phase the infection
is in.

\ChCorecap{} set up the intracellular process for such a pathogen, a
birth--death--catastrophe process with rates $\lambda$, $\mu$, $\delta$;
\ChDist{} completed its distribution theory. Both stop at the cell boundary.
This chapter carries the single-cell theory across it, and asks what
classical BMVR becomes once the cell inside it has been given
a mechanism.

\subsection{Four questions, and their answers in advance}

\begin{enumerate}
\item What replaces the release term $p\Icell$ when release is terminal and
batched?
\item In what precise sense is the classical model still right, and how
wrong is it elsewhere?
\item Does bursting help or hurt a pathogen, and at what cost?
\item How much of any of this depends on the birth--death--catastrophe
process in particular?
\end{enumerate}

The answers, in brief. What replaces $p\Icell$ is a renewal system:
infection events start cohorts, cohorts age, and two kernels weight them ---
the probability $S(\alpha)$ that a cell of age $\alpha$ is still productive,
and the expected release flux $g(\alpha)$ at that age. Both are already in
closed form from the single-cell theory, so the population model is forced
by the intracellular process rather than fitted alongside it.

The classical model is recovered exactly, but one growth rate at a time. For
each $r$ there is a unique constant pair
$\bigl(p_{\rm eff}(r),d_{\Icell,{\rm eff}}(r)\bigr)$ that reproduces the
renewal system's exponential solution exactly; what no single pair reproduces
is a transient, which samples a continuum of growth rates at once. Across
that map the release rate moves by two orders of magnitude. At
$(\lambda,\mu,\delta)=(1,0,0.1)$ the same cell presents
$p_{\rm eff}=4.59$ in a stationary infection and
$p_{\rm eff}\to\delta=0.1$ in a fast-growing one. No fit of
$(p,d_{\Icell})$ recovers the three intracellular rates: the map from
$(\lambda,\mu,\delta)$ to a constant pair has a kernel, and a third
observable is required before a population fit can test the intracellular
model.

Bursting does both. At fixed total output per cell it leaves $R_0$ exactly
where it was, lowers the establishment extinction probability when
$1/\delta<1/\lambda+1/\mu$, and slows deterministic invasion in every regime
examined. Threshold, establishment risk and invasion speed are three
different questions, and bursting answers them differently.

Almost none of it depends on the birth--death--catastrophe process. The
population equations use that process only through the pair $(S,g)$ it
supplies, so any single-cell mechanism supplying such a pair induces the
same equations, the same effective-parameter map and the same reading of
$R_0$ and $r$: bursting and budding are two choices of one object, not two
theories. What does depend on the particular process is everything resting
on its closed forms --- the flooding criterion above, the geometric burst
law, the endpoints of the map. \Cref{p:sec:reach} draws that line.

\subsection{Standing assumptions}
\label{p:sec:assumptions}

These hold throughout unless stated otherwise.

\begin{description}[leftmargin=2.2em,style=sameline]
\item[(A1)] Target cells $T$ are held constant: no depletion, no
saturation, no immune response.
\item[(A2)] Virions are not consumed on infection --- the \emph{continuum}
convention. This is why the deterministic threshold is $R_0=\gamma
TV_\infty/c$ while the discrete branching mean of \cref{p:sec:flooding} is
$m=qV_\infty$ with $q=\gamma T/(\gamma T+c)$; the two agree when $\gamma
T\ll c$.
\item[(A3)] The $\Icell_0$ cells present at $t=0$ all have age exactly
zero. A non-trivial initial age distribution adds a term to
\cref{p:eq:Icell}.
\item[(A4)] Infected cells are independent, and the population layer
assumes single infection: MOI-$k$ kernels are available
(\cref{p:app:quoted}) but are not used.
\item[(A5)] There is no spatial structure.
\end{description}

\subsection{Notation}
\label{p:sec:quoted-results}

Throughout, $\lambda$, $\mu$ and $\delta$ are the intracellular birth,
death and catastrophe rates per replicative unit. The roots $a,b$, the
derived constants $A,B,\theta,\kappa$, and the single-cell functions
$I,D,\Ifix,J,K,V$ are those of \ChCore{}, quoted in \cref{p:app:quoted}.
One naming decision is made now and kept: the burst-time density, written
$\varphi(t)$ in \ChDist{}, is written $\delta J$ here, because $\varphi$ is
reserved for the partial-release fraction of \cref{p:sec:partial}. A second
follows it: the survival kernel $S$ of \cref{p:sec:renewal} is the
non-fixation probability $\Ifix$, written $S$ where it acts as a kernel and
$\Ifix$ where the closed form is meant. Lifetime yield $V_\infty$ and mean
productive lifetime $\mean{T_{\rm prod}}$ are used from the first
calculation onward; their closed forms sit with the rest in
\cref{p:app:quoted}. \Cref{p:tab:notation} fixes the population symbols.

\begin{table}[htbp]
\centering
\renewcommand{\arraystretch}{1.12}
\small
\begin{tabular}{@{}lp{0.42\textwidth}p{0.30\textwidth}@{}}
\toprule
Symbol & Meaning & Notes \\
\midrule
$\alpha$ & cell age & argument of every kernel \\
$T$, $\gamma$, $c$ & target cells, infection rate, virion clearance & $T$ held fixed \\
$\Icell,\Vfree$ & infected cells, free particles & population counts \\
$i(t)$ & incidence, $\gamma T\Vfree$ & \\
$d_{\Icell}$ & infected-cell removal rate & phenomenological \\
$p$, $p_{\rm eff}(r)$ & constant, effective release rate & \\
$S(\alpha)$, $g(\alpha)$ & survival kernel $\Ifix(\alpha)$, release kernel $\delta K(\alpha)$ & \\
$\mathcal A(\alpha)$ & generation kernel & \cref{p:eq:genkernel} \\
$r$ & epidemic growth rate & \\
$q$, $m$ & infection success $\gamma T/(\gamma T+c)$, branching mean $qV_\infty$ & \\
$L$ & flooding parameter $a(1-b)$ & \\
$\nu$, $\varepsilon$, $Q$ & production, export rates, total stock & \cref{p:sec:odecases} \\
$\bar x$ & mean load per productive cell & $Q/\Icell$ \\
$\omega$, $d_E$, $E$ & eclipse conversion, eclipse death, eclipse class & \\
$\varphi$ & partial-release fraction & \cref{p:sec:partial} \\
$\psi$, $\zeta$, $\eta$ & release intensity, its decay and jump sizes & \cref{p:app:catalogue} \\
\bottomrule
\end{tabular}
\caption{Population notation of this chapter. Single-cell symbols are those
of \cref{p:app:quoted}.}
\label{p:tab:notation}
\end{table}

\FloatBarrier
